% RESUMO

\begin{resumo}[RESUMO (OBRIGATÓRIO)]
\begin{Spacing}{1.5}
      

Apresente de forma concisa os pontos relevantes de seu trabalho, de forma que qualquer pessoa consiga ter uma visão rápida e clara de seu conteúdo e conclusões. O resumo é formado por apenas um parágrafo, sem tópicos, contendo, preferencialmente, entre 150 e 500 palavras. Convém evitar símbolos, fórmulas, equações, contrações, reduções, entre outros, que não sejam de uso corrente e absolutamente necessários. Verbo na terceira pessoa. Nas palavras-chave indique de 3 a 6 termos que representem o conteúdo do trabalho, preferencialmente escolhidos em vocabulário controlado. NBR 6028 (Resumo).
\end{Spacing}
\vspace{1cm}
\textbf{Palavras-chave}: texto, texto, texto, texto, texto, texto.

\end{resumo}

% OBSERVAÇÕES:

% Escolha de 3 a 5 palavras ou termos que descrevam bem o seu trabalho .
% As palavras-chave são separadas por vírgula.

